% !TeX program = xelatex
% UTF-8; standalone XeLaTeX document or input into a document using xepersian.
% Complete transcription, in handwritten page order, of Notebook/coma.
% pic01.jpg = 20261012_211554878.jpg
% pic02.jpg = 20261012_211605543.jpg
% pic03.jpg = 20261012_211610754.jpg
% pic04.jpg = 20261012_211616743.jpg
% pic05.jpg = 20261012_211622455.jpg
% pic06.jpg = 20261012_211630889.jpg
% pic07.jpg = 20261012_211637841.jpg
% Source medical statements, numbers, and abbreviations are not corrected.
\ifdefined\settextfont
  \def\ComaStandalone{0}
\else
  \def\ComaStandalone{1}
  \documentclass[11pt,a4paper]{article}
  \usepackage[margin=2cm]{geometry}
  \usepackage{amsmath,amssymb,array,longtable,xcolor,enumitem}
  \usepackage{xepersian}
  \IfFontExistsTF{Amiri}{\settextfont{Amiri}}{\settextfont{FreeSerif}}
  \setlatintextfont{FreeSerif}
  \setlist[itemize]{itemsep=3pt,topsep=4pt,parsep=0pt}
\fi
\providecommand{\en}[1]{\lr{#1}}
\newcommand{\comaarr}{\ensuremath{\leftarrow}}
\newcommand{\comahl}[1]{\ifdefined\colorbox\colorbox{yellow!35}{\strut #1}\else\textbf{#1}\fi}
\newcommand{\comapage}[2]{%
  \clearpage\section*{۳ ـ کما و اختلالات هوشیاری — صفحهٔ #1}
  {\footnotesize\noindent\en{Subject:} ۳ ـ کما و اختلالات هوشیاری\par
  \noindent\en{Date: \quad / \quad /}\hfill
  \en{Sat.\quad Sun.\quad Mon.\quad Tues.\quad Thu.\quad Wed.\quad Fri.}\par
  \noindent\en{#2}\par}\medskip}
\ifnum\ComaStandalone=1 \begin{document}\fi

\comapage{۱}{pic01.jpg}
\subsection*{\en{Vegetative State (VS)}}
\begin{itemize}
  \item \comahl{۱ تا ۲ هفته پس از کمای عمیق}، برخی از اعمال مرتبط با دیانسفال و \comahl{ساقهٔ مغز} بازی می‌گردد، در حالی که عملکرد کورتکس مختل است.
  \item اگر برای \comahl{حداقل ۱ ماه} ادامه یابد \comaarr{} \comahl{\en{PVS}}.
\end{itemize}
\subsection*{تظاهرات بالینی \en{PVS}}
\begin{itemize}
  \item \comahl{باز شدن} خودبه‌خودی یا به دنبال تحریک دردناک \comahl{چشم‌ها}.
  \item حرکات \comahl{غیرارادی چشم} از یک طرف به طرف دیگر.
  \item حرکات خودبه‌خودی و \comahl{رفلکس‌های ساقهٔ مغز و نخاع وجود دارد}.
  \begin{itemize}
    \item بلع، مکیدن، جویدن، رفلکس قرنیه، ناله، \en{DTR}.
  \end{itemize}
  \item وجود \comahl{چرخهٔ خواب و بیداری}.
  \item عملکرد قلبی ـ ریوی طبیعی است.
  \item \comahl{کنترل اسفنکتر ندارد}.
\end{itemize}

\subsection*{سندرم \en{Locked-in}}
\begin{itemize}
  \item صدمات \comahl{دوطرفهٔ مسیرهای پیرامیدال در قدام پونز}.
  \item علل \comaarr{} انسداد شریان بازیلار؛ یا \comahl{میلینولیز مرکزی پونز (شایع‌ترین)}.
  \begin{itemize}
    \item گیلن‌باره \comaarr{} میاستنی گراویس \comaarr{} فلج دوره‌ای.
  \end{itemize}
  \item \comahl{بیمار بیدار و هوشیار است} / فاقد صحبت و حرکت در برابر تحریکات است.
  \item \comahl{حرکات عمودی و گاهی افقی چشم‌ها}، پلک زدن.
  \item بیماران کوادری‌پلژیک با \comahl{بابنسکی \en{upright}} دوطرفه هستند.
  \item \en{DTR} اندکی افزایش‌یافته.
  \item اصلاح سریع هایپوناترمی \comaarr{} میلینولیز مرکزی پونز.
\end{itemize}

\subsection*{مرگ مغزی}
معیارها:
\begin{enumerate}
  \item کما، عدم پاسخ به تحریکات وارده به بالای سوراخ مگنوم.
  \item تست آپنه مثبت \comaarr{} ۱۰ تا ۲۰ دقیقه از دستگاه جدا؛ \en{$PCO_2$} به ۵۰ تا ۶۰؛ به \en{$O_2$} ۱۰۰\% می‌دهیم \comaarr{} فعالیت تنفسی شروع نمی‌شود.
\end{enumerate}

\comapage{۲}{pic02.jpg}
\subsection*{مرگ مغزی — ادامهٔ معیارها}
\begin{enumerate}[start=3]
  \item فقدان رفلکس‌های سفالیک \comaarr{} مردمک، اکولوسفالیک (\en{Doll's eye})، قرنیه، اکولووستیبولار (کالریک)، \en{Gag}، بلع، مکیدن، وضعیت اکستانسور.
  \item \en{DTR}، پلنتار، فلکشن اندام \comaarr{} ممکن است وجود داشته باشد.
  \item \en{$34^{\circ}\mathrm{C}<T$} باشد.
  \item جریان خون سیستمیک ممکن است طبیعی باشد.
  \item علت باید اختلال ساختمانی مغز یا اختلال متابولیک غیرقابل برگشت باشد.
  \item \comahl{اگر احتمال مسمومیت با اتانول، خواب‌آورها، داروهای بیهوشی،}
  \comahl{داروهای فلج‌کننده مطرح باشد} \comaarr{} در تشخیص \en{BD} باید تأمل نمود.
  \item در بالغین، \comahl{حداقل ۶ ساعت} از فقدان عملکرد مغز گذشته باشد، با علت ساختمانی شناخته‌شده.
  \begin{itemize}
    \item در سایرین حداقل \comahl{۲۴ ساعت}، و از نظر دارو هم منفی.
  \end{itemize}
  \item \comahl{در نوزاد زیر ۷ روز تشخیص مرگ مغزی نگذاریم}.
  \begin{itemize}
    \item ۷ روز تا ۲ ماه: حداقل ۴۸ ساعت.
    \item ۲ ماه تا ۱ سال: حداقل ۲۴ ساعت.
    \item ۱ سال تا ۵ سال: حداقل ۱۲ ساعت (آسیب آنوکس ایسکمیک ۲۴ ساعت).
  \end{itemize}
\end{enumerate}
\subsection*{یافته‌های پاراکلینیک در \en{BD}}
\begin{itemize}
  \item \en{EEG}: در یک \comahl{دورهٔ ۳۰ دقیقه‌ای} ایزوالکتریک باشد.
  \item پتانسیل برانگیختهٔ ساقهٔ مغز (\comahl{\en{BAER}}): هیچ عکس‌العملی ایجاد نمی‌شود.
  \item در مطالعات تصویربرداری مثل \en{MRA}، \en{TCD} و آنژیوگرافی، \comahl{جریان خون مغزی} وجود ندارد.
\end{itemize}

\comapage{۳}{pic03.jpg}
\subsection*{\comahl{\en{GCS}}}
\subsubsection*{باز شدن چشم‌ها}
\begin{center}
\begin{tabular}{|p{10cm}|c|}
\hline
باز شدن خودبه‌خودی & ۴ \\ \hline
با صحبت & ۳ \\ \hline
با تحریک دردناک & ۲ \\ \hline
باز نمی‌کند & ۱ \\ \hline
\end{tabular}
\end{center}
\subsubsection*{پاسخ کلامی}
\begin{center}
\begin{tabular}{|p{10cm}|c|}
\hline
صحبت با مفهوم & ۵ \\ \hline
محاورهٔ نامفهوم & ۴ \\ \hline
کلمات نامربوط & ۳ \\ \hline
صداهای غیرقابل فهم & ۲ \\ \hline
سکوت & ۱ \\ \hline
\end{tabular}
\end{center}
\subsubsection*{پاسخ حرکتی}
\begin{center}
\begin{tabular}{|p{10cm}|c|}
\hline
انجام کامل دستورات & ۶ \\ \hline
تعیین محل درد & ۵ \\ \hline
عقب کشیدن در برابر درد & ۴ \\ \hline
فلکشن غیرطبیعی اندام فوقانی & ۳ \\ \hline
اکستنشن اندام فوقانی & ۲ \\ \hline
عدم واکنش & ۱ \\ \hline
\end{tabular}
\end{center}

\subsection*{علل هایپوترمی منجر به کما (\en{$T\downarrow$})}
\begin{itemize}
  \item مسمومیت با الکل (انسفالوپاتی ورنیکه).
  \item مسمومیت با باربیتورات‌ها.
  \item غرق شدن / مواجهه با هوای سرد.
  \item میگزدم ناشی از هایپوتیروئیدی.
  \item اختلال جریان خون محیطی.
\end{itemize}
\subsection*{علل هایپرترمی منجر به کما (\en{$T\uparrow$})}
\begin{itemize}
  \item مسمومیت با آنتی‌کولینرژیک‌ها.
  \item عفونت‌های سیستمیک.
  \item مننژیت / انسفالیت.
  \item گرمازدگی \comaarr{} ضایعهٔ مرکز تنظیم‌کنندهٔ دما در مغز.
\end{itemize}

\comapage{۴}{pic04.jpg}
\begin{itemize}
  \item در مسمومیت‌های دارویی و اختلالات متابولیک منجر به کما \comaarr{} \comahl{مردمک‌ها اغلب طبیعی}.
  \item مردمک ثابت و دیلاتهٔ یک‌طرفه:
  \begin{itemize}
    \item سایز مردمک \en{$>5.5\,\mathrm{mm}$} / به نور واکنش نمی‌دهد / فشار به عصب \en{III} (اکولوموتور)، به دلیل تودهٔ فضاگیر همان طرف (هرنی مغزی).
  \end{itemize}
  \item مردمک ثابت و دیلاتهٔ دوطرفه:
  \begin{itemize}
    \item فتق ترنس‌تنتوریال.
    \item مسمومیت با \en{TCA}؛ مسمومیت با آتروپین.
  \end{itemize}
  \item مردمک میوتیک \comaarr{} مسمومیت باربیتورات‌ها / سندرم هورنر.
  \item مردمک \en{pinpoint} \comaarr{} خونریزی پونز / مسمومیت اپیوئید.
  \item مردمک \en{Hippus} \comaarr{} تغییر متناوب سایز در پاسخ به نور.
  \begin{itemize}
    \item کاراکتریستیک انسفالوپاتی متابولیک.
  \end{itemize}
  \item رفلکس \en{Cilio-Spinal} \comaarr{} فشار به لترال گردن \comaarr{} گشاد شدن مردمک همان سمت.
  \begin{itemize}
    \item در \comahl{ضایعات ساقهٔ مغز} از بین می‌رود.
  \end{itemize}
  \item ضایعات \en{Midbrain} \comaarr{} مردمک متوسط و \en{fixe}.
  \item ضایعات دیانسفال \comaarr{} مردمک کوچک و واکنشی.
  \item ضایعات پره‌تکتال \comaarr{} مردمک دیلاته و \en{fixe}.
\end{itemize}

\subsection*{حرکات چشم‌ها در کما}
\begin{itemize}
  \item کمای خفیف ناشی از اختلالات متابولیک \comaarr{} چشم‌ها بدون هدف خاصی از یک طرف به طرف دیگر حرکت می‌کند.
  \item ضایعات وسیع یک نیمکره \comaarr{} انحراف چشم‌ها به سمت ضایعه یا به سمت مخالف همی‌پارزی.
  \item ضایعات یک‌طرفهٔ پونز \comaarr{} انحراف چشم‌ها به سمت مخالف ضایعه یا به طرف همی‌پارزی.
\end{itemize}

\comapage{۵}{pic05.jpg}
\subsection*{رفلکس اکولوسفالیک (\en{Doll's eye})}
\begin{itemize}
  \item مرکز کنترل حرکات افقی چشم ساقهٔ مغز (\en{brain stem}) است.
  \item مثبت \comaarr{} چشم‌ها به سمت مخالف حرکت کند \comaarr{} سالم بودن ساقهٔ مغز و وجود ضایعات سوپراتنتوریال.
\end{itemize}
\begin{itemize}
  \item خونریزی پونز \comaarr{} شروع ناگهانی کما، مردمک \en{pinpoint}، چشم حرکت افقی ندارد ولی حرکت عمودی \en{NL}، کوادری‌پارزی، وضعیت دسربره، هایپرترمی (\en{$39.5<T$}).
  \item در مسمومیت با اپیوئید، حرکات افقی چشم‌ها \en{NL} است.
\end{itemize}

\subsection*{تست کالریک}
بررسی عملکرد سیستم \comahl{وستیبولار}.
\begin{itemize}
  \item رفلکس اکولووستیبولار \comaarr{} در صورت عملکرد طبیعی ساقهٔ مغز، به افراد کمک می‌کند هنگام حرکت سر چشم‌هایشان را بر یک هدف بینایی ثابت نگه دارند.
  \item تست کالریک فرد سالم پس از تزریق آب سرد به گوش:
  \begin{itemize}
    \item فاز کند نیستاگموس: تحریک سیستم وستیبولار همان گوش و چشم‌ها به سمت تحریک می‌رود.
    \item فاز تند نیستاگموس: چشم‌ها سریعاً به سمت گوش مقابل حرکت می‌کند.
  \end{itemize}
  \item تست کالریک و بیمار کمایی:
  \begin{itemize}
    \item تست کالریک مثبت (رفلکس مثبت) \comaarr{} صرفاً حرکت چشم‌ها به سمت گوش تحریک‌شده:
    \begin{itemize}
      \item ساقهٔ مغز سالم است ولی کورتکس مشکل دارد.
    \end{itemize}
    \item تست کالریک منفی (رفلکس منفی) \comaarr{} هیچ حرکتی در چشم‌ها دیده نمی‌شود.
    \begin{itemize}
      \item ساقهٔ مغز درگیر است.
    \end{itemize}
    \item با آب گرم نتایج برعکس خواهد بود.
  \end{itemize}
\end{itemize}

\comapage{۶}{pic06.jpg}
\begin{itemize}
  \item مسمومیت با ضدافسردگی‌ها و ضدصرع‌ها در موارد خفیف‌تر رفلکس اکولوسفالیک و در موارد شدیدتر رفلکس اکولووستیبولار را منفی می‌کند.
\end{itemize}
\subsection*{الگوهای تنفسی در بیماران کمایی}
\begin{itemize}
  \item تنفس آتاکسیک یا بیوت \comaarr{} ضایعات مدولا.
  \begin{itemize}
    \item نامنظم، بدون ریتم خاص، تنفس‌ها از نظر ریتم و عمق فرق دارند.
    \item تنفس‌ها بدتر شده \comaarr{} منجر به آپنهٔ کامل تنفسی.
  \end{itemize}
  \item تنفس شین ـ استوکس \comaarr{} پیش‌آگهی بدی ندارد.
  \begin{itemize}
    \item به دلیل جدا شدن مراکز تنفسی ساقهٔ مغز از کورتکس و افزایش حساسیت آن‌ها نسبت به \en{$CO_2$}.
    \item دوره‌های متناوب و منظم به صورت افزایش و سپس کاهش تنفس و سرانجام آپنهٔ تنفسی.
    \item ضایعات دوطرفهٔ نیمکره‌ها، به‌خصوص نواحی عمقی و دیانسفال.
    \item هماتوم ساب‌دورال دوطرفه.
    \item مسمومیت دارویی یا اختلالات متابولیک.
  \end{itemize}
  \item تنفس تند و عمیق مرکزی (سنترال نورولوژیک هایپرونتیلاسیون) (\en{CNH}):
  \begin{itemize}
    \item تعداد و عمق تنفس افزایش می‌یابد \comaarr{} منجر به آلکالوز تنفسی می‌شود.
    \item نام دیگر آن تنفس کوسمال است.
    \item ممکن است هیچ ارتباط پایداری بین این تنفس و مکان ضایعه وجود نداشته باشد.
    \item ضایعات مناطق تگمان مدبرین، فوقانی تگمنتوم پونز.
    \item لنفوم اولیهٔ مغز بدون درگیری ساقهٔ مغز.
    \item کمای متابولیک.
  \end{itemize}
\end{itemize}

\comapage{۷}{pic07.jpg}
\begin{itemize}
  \item تنفس آپنوستیک \comaarr{} وقفهٔ ۲ تا ۳ ثانیه‌ای بعد از هر دم.
  \begin{itemize}
    \item ضایعات پایین پونز به علت انسداد شریان بازیلار.
  \end{itemize}
\end{itemize}

\begin{itemize}
  \item وضعیت دکورتیکه \comaarr{} ضایعات نیمکرهٔ مغزی / کپسول داخلی / تالاموس.
  \item وضعیت دسربره \comaarr{} فشار به مدبرین به دلیل ضایعات فضاگیر نیمکره‌های مغزی یا ضایعات \en{posterior fossa}.
  \begin{itemize}
    \item کمای متابولیک (آنوکسی، هیپوگلیسمی، کمای کبدی، مسمومیت شدید).
  \end{itemize}
\end{itemize}

\subsection*{اقدامات اورژانسی در کما}
\begin{itemize}
  \item مرحلهٔ اول: \en{ABC} / نمونهٔ خون (قند، الکترولیت، کبد، کلیه، \en{CBC}، \en{PT}، \en{PTT}).
  \begin{itemize}
    \item تزریق \en{IV} ۲۵ گرم دکستروز، ۱۰۰ میلی‌گرم تیامین، \en{0.4--2} میلی‌گرم نالوکسان / \en{ABG} / مهار تشنج در صورت وقوع.
  \end{itemize}
  \item مرحلهٔ دوم: اگر علائم تحریک مننژ وجود داشت \comaarr{} \en{LP}.
  \begin{itemize}
    \item معاینهٔ فیزیکی و نورولوژیک دقیق.
    \item انجام \en{CT} سریع \comaarr{} وجود یافته‌ای به نفع یک ضایعهٔ ساختمانی.
    \begin{itemize}
      \item تروما، به نفع خونریزی \en{SAH}.
      \item قبل از انجام \en{LP}.
    \end{itemize}
  \end{itemize}
  \item مرحلهٔ سوم: \en{ECG}، اصلاح هایپو و هایپرترمی، اصلاح اختلالات اسید و باز و الکترولیت، \en{CXR}، آزمایش ادرار و خون از نظر مسمومیت، \en{EEG}.
\end{itemize}

% Handwritten source page numbers: 1--7. Printed Date and weekday labels
% are reproduced in each page header; the date and weekday fields are blank.
\ifnum\ComaStandalone=1 \expandafter\enddocument\fi
